\subsection{Silver-Supervised Adaptation}
\label{sec:silver-results}

Silver supervision improved Qwen's typed exact-span micro-F1 under the fixed inference protocol. The adapted mean was $0.540\pm0.035$, compared with 0.361 without an adapter (Figure~\ref{fig:model-comparison}b; Table~\ref{tab:gold150-main}, panel A). The mean gain over the three training seeds was 0.179, with a paired 95\% bootstrap interval of $[0.111, 0.242]$; the intervals for all three individual seeds also lay above zero (Table~\ref{tab:core-paired}). Mean scores also increased under relaxed and boundary-only matching, showing gains across the three matching criteria.

The adapters produced fewer responses with rejected span candidates than the unadapted model. These end-to-end scores reflect both span extraction and output validity. The reproduction materials include \href{https://github.com/AlfredJamesLi/chinese-skillspan-benchmark/blob/1d1c24e20ee049fc768034b5baf58fa882e08e2c/reproduction/process_archive_20260924/secondary_review/tables/qwen_per_seed.csv}{per-seed Qwen scores} and parser outcomes.

The silver-trained JobBERT-zh + CRF baseline achieved mean typed exact-span micro-F1 of $0.554\pm0.005$ (Figure~\ref{fig:model-comparison}b; Table~\ref{tab:gold150-main}, panel B). Its mean exceeded Qwen + LoRA (0.540), while both remained below prompted GPT-5.6 Terra (0.667). These rankings describe the evaluated configurations, whose training and inference protocols differ across model families.
